% GPT-6 Astra Pro; effort not exposed; /root/replacement; 22 November 2026.
\section[HD-F1-R25: target surface defects]{HD-F1-R25: the common-image target must have a surface defect in one of its rank-one factors}
\label{sec:followup-rm-common-image-target-surface}

Retain exactly R23's actual common-image alternative:
an RM fourfold of type $(2,2)$, the actual determinant
twist of the specified finite coherent point action,
and a gluable weakly semiregular adjacent object
\[
 A\longrightarrow E\longrightarrow T[-1]
       \xrightarrow{\delta[-1]}A[1],\qquad
 \operatorname{ch}(E)=v_q=(1-qA_1)(1-qA_2),\quad q>0.
 \tag{F1R25.1}
\]
Here $A,T$ are equal-slope semistable TF sheaves of
ranks $(3,2)$, their common slope is zero, and every
invariant map $A\to T$ factors through the specified
saturated rank-one subsheaf $I\subset T$. Write
\[
 Q=T/I,\qquad L=I^{**},\qquad M=Q^{**},\qquad
 Z_I=L/I,\qquad Z_Q=M/Q.
 \tag{F1R25.2}
\]
Thus $Q$ is TF of rank one, $L,M$ are actual line
bundles, and all displayed objects and maps carry
the induced lifts for that same action. We prove
\[
                [Z_I]_2+[Z_Q]_2\ne0.
 \tag{F1R25.3}
\]
The subscript denotes the actual effective
codimension-two support cycle, with generic
lengths. In particular, at least one of the two
rank-one factors of the target has a surface
defect. This holds for both $h=2$ and $h=3$
from R23. Together with R23 it requires a
source or evaluation surface defect and a
target-factor surface defect on the same
candidate; it retains their actual possible
cancellation.

\paragraph{1. An exact lift of the target across smaller-dimensional defects.}
Assume for contradiction that both $Z_I$ and $Z_Q$
have support of codimension at least three.
The actual equivariant pushout along $I\hookrightarrow L$
of $0\to I\to T\to Q\to0$ is
\[
 \begin{gathered}
 0\longrightarrow L\longrightarrow T_1
                     \longrightarrow Q\longrightarrow0,\\
 0\longrightarrow T\longrightarrow T_1
                     \longrightarrow Z_I\longrightarrow0.
 \end{gathered}
 \tag{F1R25.4}
\]
For any sheaf $Z$ with this codimension of support
and any line bundle $B$, the sheaf groups
$\mathcal Ext^j(Z,B)$ vanish for $j<3$.
The local-to-global Ext spectral sequence consequently
gives
\[
                         \operatorname{Ext}^j(Z,B)=0
                                  \quad(j=0,1,2).
 \tag{F1R25.5}
\]
This is the exact low-degree vanishing in
Huybrechts--Lehn, Proposition~1.1.6(i), after the
actual line twist. Taking finite-group invariants
preserves it.

Apply $\operatorname{Hom}(-,L)$ to
$0\to Q\to M\to Z_Q\to0$. The two neighboring
terms in degrees one and two vanish by (F1R25.5),
so the actual restriction map is an isomorphism
\[
 \operatorname{Ext}^1_G(M,L)
          \xrightarrow{\ \sim\ }\operatorname{Ext}^1_G(Q,L).
 \tag{F1R25.6}
\]
Here $\operatorname{Ext}^1_G=\operatorname{Ext}^1{}^G$:
the induced action on Hom complexes is a genuine
finite-group action for the given coherent lifts,
and the averaging projector makes its invariants
exact. In particular (F1R25.6) concerns extension
classes of those actual lifts. The class of the
first sequence in (F1R25.4) therefore gives an
actual equivariant extension
\[
 0\longrightarrow L\xrightarrow{i}\mathcal V
                 \xrightarrow{p}M\longrightarrow0,
       \qquad e:M\longrightarrow L[1].
 \tag{F1R25.7}
\]
Since its two endpoints are locally free,
$\mathcal V$ is locally free. Its pullback along
$Q\hookrightarrow M$ is the first sequence in
(F1R25.4), including its actual identifications.
Thus there is an equivariant monomorphism
$j:T\hookrightarrow\mathcal V$, with quotient $Z$
fitting into
\[
                 0\longrightarrow Z_I\longrightarrow Z
                          \longrightarrow Z_Q\longrightarrow0.
 \tag{F1R25.8}
\]
It follows from (F1R25.5), with the degree-two
predecessor retained, that
\[
 j^*:\operatorname{Ext}^2_G(\mathcal V,M)
                 \lhook\joinrel\longrightarrow
                     \operatorname{Ext}^2_G(T,M).
 \tag{F1R25.9}
\]
Indeed the predecessor in the Hom sequence is
$\operatorname{Ext}^2_G(Z,M)=0$. This is the
codimension-three application of that exact
sequence; no vanishing in codimension two is used.

\paragraph{2. The actual endpoint relation becomes a fixed cup-product image.}
Put $K=\ker c^2_{v_q}$. The weak identity-sector
comparison in R05 and the actual endpoint Hom
sequence in R19 give, for every $k\in K$, an
invariant $f_k:A\to T$ with
\[
                    \operatorname{ev}_T^2(k)=f_k[2]\delta.
 \tag{F1R25.10}
\]
Every $f_k$ has image in $I$, and $j(I)\subset i(L)$,
so $p j f_k=0$. Naturality for the actual map
$p j:T\to M$ now gives
\[
 \bigl(\operatorname{ev}_M^2(k)p\bigr)j
       =(p j)[2]\operatorname{ev}_T^2(k)
       =(p j f_k)[2]\delta=0.
 \tag{F1R25.11}
\]
By (F1R25.9), $\operatorname{ev}_M^2(k)p=0$ in
$\operatorname{Ext}^2_G(\mathcal V,M)$.
The Hom sequence of (F1R25.7) is
\[
 \operatorname{Ext}^1_G(L,M)
   \xrightarrow{\,v\mapsto v[1]e\,}
 \operatorname{Ext}^2_G(M,M)
   \xrightarrow{\,p^*\,}
 \operatorname{Ext}^2_G(\mathcal V,M).
 \tag{F1R25.12}
\]
Forgetting equivariance in this exact membership
and using the actual line-tensor identifications,
we obtain
\[
 \operatorname{ev}_M^2(K)\ \subset\
 \mathcal C_e:=
 \operatorname{im}\left[
       H^1(X,D^{-1})\xrightarrow{\,v\mapsto v[1]e\,}
       H^2(X,\mathcal O_X)\right],
 \qquad
 D=L\otimes M^{-1},\quad e\in H^1(X,D).
 \tag{F1R25.13}
\]
This is the Yoneda cup product for the fixed actual
extension (F1R25.7). Enlarging the invariant source
to its ordinary cohomology in (F1R25.13) only
enlarges its image. No hypothesis on ordinary
cross-Hom groups is imposed.

\paragraph{3. Every such line-bundle cup image has a common exterior factor.}
We prove the needed assertion on an arbitrary
complex abelian fourfold: for every actual line
bundle $D$ and every fixed $e\in H^1(D)$, there
is a nonzero $u\in P:=H^1(\mathcal O_X)$ with
\[
                              \mathcal C_e\subset u\wedge P
                                   \subset\bigwedge^2P.
 \tag{F1R25.14}
\]
The zero image satisfies this for any nonzero $u$.
Here are the explicit quotient, pullback and cup maps.

Write $X=\mathbb V/\Lambda$. By Appell--Humbert,
$D$ is represented by a Hermitian form $H_D$ and
a unitary semicharacter $\chi$ on $\Lambda$, with
$\operatorname{Im}H_D$ integral on $\Lambda^2$.
With $H_D$ complex linear in the first variable
its factors of automorphy are
\[
 j_\lambda(z)=\chi(\lambda)
       \exp\bigl(\pi H_D(z,\lambda)+\tfrac{\pi}{2}
                               H_D(\lambda,\lambda)\bigr).
 \tag{F1R25.15}
\]
Let $\mathbb V_0=\operatorname{rad}H_D$ and
$\Lambda_0=\Lambda\cap\mathbb V_0$. The real kernel
of the integral alternating form
$\operatorname{Im}H_D$ is rational; its type $(1,1)$
makes it complex, and it equals this Hermitian
radical. Thus $\Lambda_0$ is a full lattice in
$\mathbb V_0$ and $X_0=\mathbb V_0/\Lambda_0$ is
an abelian subvariety.

Choose an actual ample polarization with positive
Hermitian form $H_0$, and let
$\mathbb V_1=\mathbb V_0^{\perp H_0}$,
$\Lambda_1=\Lambda\cap\mathbb V_1$.
This complement is complex and rational: it is
also the real orthogonal complement for the
integral alternating form $\operatorname{Im}H_0$,
and that form is nondegenerate on $\mathbb V_0$.
The actual addition map
\[
 a:X_0\times X_1\longrightarrow X,\qquad
       X_1=\mathbb V_1/\Lambda_1,\qquad (x_0,x_1)\longmapsto x_0+x_1
 \tag{F1R25.16}
\]
is a finite etale isogeny: its differential is the
direct-sum isomorphism
$\mathbb V_0\oplus\mathbb V_1\to\mathbb V$, and
$\Lambda_0+\Lambda_1$ has finite index in $\Lambda$.
Restriction of (F1R25.15) gives the actual isomorphism
\[
                        a^*D=D_0\boxtimes D_1,
 \tag{F1R25.17}
\]
where $D_0$ is the flat unitary line of character
$\chi|_{\Lambda_0}$ and $D_1$ is nondegenerate on
$X_1$. There are no mixed factors, since $H_D$
vanishes whenever an argument lies in $\mathbb V_0$.

A nontrivial flat unitary line on a complex torus
has zero cohomology in every degree. Explicitly,
lift its Dolbeault complex to the covering vector
space and use the unitary Fourier modes indexed by
a real character shift plus the dual lattice.
On each mode the differential is wedge by its
nonzero $(0,1)$ frequency. Contraction by the
metric dual frequency, divided by its squared norm,
is a contracting homotopy. A real frequency has
zero $(0,1)$ component only when it is zero;
the zero frequency occurs exactly for the trivial
unitary character. The shifted lattice stays
away from zero, so these contractions preserve
the rapidly decreasing Fourier coefficients.
This proves the assertion for the full smooth
Dolbeault complex.

For the finite flat map $a$, the unit
$D\to a_*a^*D$ followed by the actual trace is
multiplication by $\deg(a)$. Therefore $a^*$ is
injective on $H^j(D)$, with retraction
$\operatorname{Tr}_a/\deg(a)$.
Kunneth in (F1R25.17) shows that nonzero $e$
requires $D_0=\mathcal O_{X_0}$.
In that case $\chi|_{\Lambda_0}=1$, so
$H_D$ and $\chi$ descend respectively to
$\mathbb V/\mathbb V_0$ and $\Lambda/\Lambda_0$:
the semicharacter rule gives
$\chi(\lambda+\lambda_0)=\chi(\lambda)$.
Formula (F1R25.15) on these quotients constructs
an actual nondegenerate line $\overline D$ on
$\overline X=X/X_0$, and the quotient map
$b:X\to\overline X$ has
\[
                             D=b^*\overline D.
 \tag{F1R25.18}
\]
Thus both descent and the product pullback in
this computation are specified maps.

Put $r=\dim X_1$, and let $s$ be the index of
$D_1$. Mumford's nondegenerate vanishing and index
theorems say that its cohomology is concentrated
in degree $s$, while that of $D_1^{-1}$ is
concentrated in degree $r-s$.
The latter equality also follows directly from
Serre duality and the trivial canonical bundle.
Write $P_j=H^1(X_j,\mathcal O_{X_j})$. Kunneth
gives the actual degree-one groups
\[
 \begin{split}
 H^1(a^*D)&=H^s(D_1)\otimes\bigwedge^{1-s}P_0,\\
 H^1(a^*D^{-1})&=H^{r-s}(D_1^{-1})
                         \otimes\bigwedge^{1-r+s}P_0,
 \end{split}
 \tag{F1R25.19}
\]
with negative exterior degrees interpreted as zero.
Both groups can be nonzero only for
$s\le1$ and $r-s\le1$. Consequently $r\le2$.
Under the actual Kunneth cup map their products
have the following complete possibilities:
\[
 \begin{array}{c|c|c}
 r & \text{possible }s & \text{subspace containing the product}\\ \hline
 0 & 0 & (a^*e)\wedge P_0\\
 1 & 0,1 & P_1\wedge P_0,\quad \dim P_1=1\\
 2 & 1 & \bigwedge^2P_1,\quad \dim P_1=2 .
 \end{array}
 \tag{F1R25.20}
\]
For $r=0$, (F1R25.18) gives $D=\mathcal O_X$.
For $r=1$ every product contains a fixed nonzero
generator of $P_1$ as an exterior factor; this
holds also when $a^*e$ is a sum of Kunneth tensors.
For $r=2$ the displayed subspace is a decomposable
line and has any prescribed nonzero vector of
$P_1$ as a common exterior factor. For $r>2$ one
of the two groups in (F1R25.19) is zero.

Finally $a^*:P\to P_0\oplus P_1$ is an isomorphism:
in invariant Dolbeault forms it is the dual
antiholomorphic differential of (F1R25.16).
Its degree-two map is its exterior square.
The square comparing the original cup product
with the Kunneth cup product commutes by
pullback of Dolbeault representatives. Taking
the inverse of this degree-one isomorphism
transports the common factor in (F1R25.20)
to the required $u\in P$, proving (F1R25.14).
The auxiliary isogeny has only been used to
compute this ordinary cup map; no new action
or substitute object is attached to the candidate.

\paragraph{4. The balanced line action has no common exterior factor.}
The explicit R07 quadratic relations are
\[
 U=U_1\oplus U_2,\qquad U_i=T_i\oplus P_i,\qquad
 P=P_1\oplus P_2,\qquad
 K=K_1\oplus K_2,\qquad K_i=\ker\omega_i,
 \tag{F1R25.21}
\]
where $\dim T_i=\dim P_i=2$ and $\omega_i$ is a
nondegenerate alternating form on $U_i$.
Here $P_i$ denote the RM blocks; they are
independent of the auxiliary product notation
in paragraph3.
For the actual line $M$, with $\ell=c_1(M)$,
the degree-one map and its exterior square are
\[
 \pi_\ell:U\longrightarrow P,\qquad
              \pi_\ell(v,\eta)=\eta+\iota_v\ell,
       \qquad
              \operatorname{ev}_M^2=\bigwedge^2\pi_\ell.
 \tag{F1R25.22}
\]
Indeed the full trace of a class
$z\in\operatorname{Ext}^2(M,M)=H^2(\mathcal O_X)$
is $e^\ell z$, while the actual characteristic
action on $e^\ell$ is
$e^\ell\bigwedge^2\pi_\ell$. Multiplication by
$e^\ell$ is injective, proving the last equality
with the same sign convention as R09 and R22.
In particular $\pi_\ell|_{P_i}$ is the identity.

Fix any nonzero $u\in P$ and let
$\rho_u:P\to P/\mathbf C u$ be the actual quotient.
At least one of the two spaces $\rho_u(P_i)$ has
dimension two, since $P_1\cap P_2=0$.
For that index the map
\[
                    \varphi_i=\rho_u\pi_\ell|_{U_i}
 \tag{F1R25.23}
\]
has rank at least two.
Suppose its exterior square vanished on $K_i$.
Choose two linear functionals on its image whose
pullbacks to $U_i$ are independent. Their wedge
is a nonzero decomposable alternating form on
$U_i$ vanishing on $K_i$. Since $K_i$ is the
hyperplane $\ker\omega_i$, that form is a
nonzero scalar multiple of $\omega_i$.
The former has rank two and the latter rank
four, a contradiction. Thus
\[
                    \bigwedge^2\varphi_i(K_i)\ne0.
 \tag{F1R25.24}
\]
The kernel of $\bigwedge^2\rho_u$ is exactly
$u\wedge P$, as is seen by extending $u$ to a
basis of $P$. Equations (F1R25.22)--(F1R25.24)
therefore prove
\[
        \operatorname{ev}_M^2(K)\not\subset u\wedge P
                              \quad\text{for every }0\ne u\in P.
 \tag{F1R25.25}
\]
This proof permits every actual algebraic line
$M$; its divisor class is not restricted to
$F\theta$.

\paragraph{5. Contradiction and exact surviving scope.}
The exact endpoint membership (F1R25.13) and
the cup-image calculation (F1R25.14) contradict
(F1R25.25). Hence $Z_I$ and $Z_Q$ cannot both
have codimension at least three.
Both already have codimension at least two,
and their leading support cycles are effective.
This proves (F1R25.3).

The obstruction uses the full balanced
Hochschild kernel and actual equivariant
extensions. Its only line-bundle cohomology
premise is Mumford's published vanishing and
index theorem, applied after the explicit
Appell--Humbert quotient and isogeny maps.
Mumford, \emph{Abelian Varieties}, Chapter~I,
\S2, pp.~14--22, supplies the factors of
automorphy, their Chern class and the complete
Appell--Humbert proof; Chapter~III, \S16,
supplies the nondegenerate cohomology theorem
already checked for R14. The exact codimension
vanishing and its proof are
Huybrechts--Lehn, second edition,
Proposition~1.1.6(i). The characteristic/trace
comparison is the actual Buchweitz--Flenner
map used in R07, R09 and R19.

When the surface parts in (F1R25.3) occur,
the two Ext groups removed in (F1R25.6) and
(F1R25.9) can contribute. No injectivity
or extension lift across those surface defects
has been asserted. In particular the argument
does not close the surviving surface-cancellation
branches, and it does not replace the actual
coupling $\delta$ by a numerical class.
